% ============================================================
% Section 6: Conclusion
% ============================================================

\section{Conclusion}
\label{sec:conclusion}

We presented DocVerify, a multi-task CNN for joint detection and localization of
forged identity documents, paired with a multi-objective Pareto HPO framework
that selects hyperparameter configurations by optimizing PR-AUC and Dice
simultaneously without fixing their relative importance in advance.

Our main findings are: \rev{(i) learning detection and localization jointly
benefits both tasks: a single-task model performs well only on the axis it
targets (detection-only Dice $0.002$; localization-only PR-AUC $0.728$; both
$p < 0.001$), whereas the joint model attains both; it further yields higher
Dice than equal-weight training ($p = 1.3 \times 10^{-4}$, $d = 1.01$) at no
PR-AUC cost, and distance-to-ideal selection outperforms scalar
PR-AUC/Dice-maximizing selection on the same front}; (ii) under a rigorous 10-fold
nested cross-validation protocol, DocVerify achieves PR-AUC $= 0.9921 \pm
0.0058$ and Dice $= 0.856 \pm 0.030$, demonstrating consistent generalization
across data partitions; and \rev{(iii) trained on FantasyID alone, DocVerify
attains F1@0.5 $= 0.969 \pm 0.014$ and F1$_{\text{loc}} = 0.807 \pm 0.096$ on our
internal holdout, reported alongside the DeepID 2025 Challenge for context only,
as the evaluation sets differ.}

\rev{Applying the identical MOTPE protocol to a ResNet-18 backbone yields
comparable blind-test performance (Sec.~\ref{subsec:generalization}), so the
approach is not tied to one architecture. By contrast, zero-shot evaluation on
SIDTD exposed a substantial cross-domain gap; closing it through multi-source
training is a priority for future work}, alongside forgery-artifact features
(noise fingerprints, frequency anomalies~\citep{cozzolino2020noiseprint}) and
foundation model backbones~\citep{fakeidet2025} for richer pre-training.

\section*{Acknowledgments}
This work was supported by the Spanish National Cybersecurity Institute (INCIBE,
Spain). The author thanks the Idiap Research Institute for releasing the
FantasyID dataset.
